% GENERATED FILE. Edit canonical lessons, then run npm run generate:books.
% canonical-source-hash: fnv1a64:10c2ba15e0526127
% canonical-lessons: GU-C541-sapath, GU-C541-vaydo, GU-C541-bharoso, GU-C541-chatakbari, GU-C541-prastav

\chapter{Promises and Trust}
\label{ch:gu-promises-and-trust}
\hlchaptermodality{541}

\begin{chapteropening}
\textbf{By the end of this chapter:} \emph{I can talk about oaths, promises, trust and loopholes.}

Complete the last lesson of chapter 541: I can talk about oaths, promises, trust and loopholes.
\end{chapteropening}

\section[\texorpdfstring{\gu{શપથ} (śapath)}{શપથ (śapath)}]{\gu{શપથ} (śapath) — a solemn oath}

\label{lesson:GU-C541-sapath}



\begin{quote}
Before the new one: say the Gujarati for to hold fast, then the Gujarati for an oath.
\end{quote}

\subsection*{You'll want to know: \gu{શપથ}}
\textbf{\gu{શપથ}} — \emph{śapath} — \textquotedblleft{}a solemn oath\textquotedblright{}.

\begin{quote}
\gu{જો} \gu{તમે} \gu{શપથ} \gu{લેશો} \gu{તો} \gu{તમારે} \gu{સાચું} \gu{બોલવું} \gu{પડશે}.
\end{quote}

\begin{quote}
\emph{jo tame śapath leśo to tamāre sācũ bolvũ paḍśe.}
\end{quote}

\begin{quote}
If you take the oath, you will have to tell the truth.
\end{quote}

A masculine noun.

\begin{cousinweb}
\textbf{\gu{શપથ}} (\emph{śapath}) is a tatsama, taken from Sanskrit \emph{śapatha}, \textquotedblleft{}a curse; an oath\textquotedblright{} (Monier-Williams). One takes it formally: \textbf{\gu{શપથ} \gu{લેવા}}.
\end{cousinweb}

\begin{grammarlens}[title={What you've built}]
One more word for saying what happens if.
\end{grammarlens}

\subsection*{Your turn}
\begin{itemize}
\raggedright
  \item \emph{Say it:} \emph{śapath}
  \item \emph{Say it:} \emph{śapath}, once more
  \item \emph{Say it:} \emph{jo tame śapath leśo to tamāre sācũ bolvũ paḍśe.} — If you take the oath, you will have to tell the truth.
  \item \emph{Recall:} say the Gujarati for to hold fast, then the Gujarati for an oath, then say \emph{śapath} again
\end{itemize}

\begin{tcolorbox}[breakable,colback=teal!4,colframe=teal!35!black,title={Before you move on}]
What does \gu{શપથ} mean? (\textquotedblleft{}A solemn oath\textquotedblright{}.) Say it once more.
\end{tcolorbox}
\section[\texorpdfstring{\gu{વાયદો} (vāydo)}{વાયદો (vāydo)}]{\gu{વાયદો} (vāydo) — an undertaking (to do something by a time)}

\label{lesson:GU-C541-vaydo}



\begin{quote}
Before the new one: say the Gujarati for an oath, then the Gujarati for a solemn oath.
\end{quote}

\subsection*{You'll want to know: \gu{વાયદો}}
\textbf{\gu{વાયદો}} — \emph{vāydo} — \textquotedblleft{}an undertaking (to do something by a time)\textquotedblright{}.

\begin{quote}
\gu{જો} \gu{તમે} \gu{વાયદો} \gu{કરશો} \gu{તો} \gu{તમારે} \gu{તે} \gu{પાળવો} \gu{પડશે}.
\end{quote}

\begin{quote}
\emph{jo tame vāydo karśo to tamāre te pāḷvo paḍśe.}
\end{quote}

\begin{quote}
If you make a promise, you will have to keep it.
\end{quote}

A masculine noun.

\begin{cousinweb}
is usually traced to Arabic \emph{waʿda}, \textquotedblleft{}a promise\textquotedblright{}, through Persian. \textbf{\gu{વચન}} is one's word; \textbf{\gu{વાયદો}} is often a promise with a date in it.
\end{cousinweb}

\begin{grammarlens}[title={What you've built}]
One more word for saying what happens if.
\end{grammarlens}

\subsection*{Your turn}
\begin{itemize}
\raggedright
  \item \emph{Say it:} \emph{vāydo}
  \item \emph{Say it:} \emph{vāydo}, once more
  \item \emph{Say it:} \emph{jo tame vāydo karśo to tamāre te pāḷvo paḍśe.} — If you make a promise, you will have to keep it.
  \item \emph{Recall:} say the Gujarati for an oath, then the Gujarati for a solemn oath, then say \emph{vāydo} again
\end{itemize}

\begin{tcolorbox}[breakable,colback=teal!4,colframe=teal!35!black,title={Before you move on}]
What does \gu{વાયદો} mean? (\textquotedblleft{}An undertaking (to do something by a time)\textquotedblright{}.) Say it once more.
\end{tcolorbox}
\section[\texorpdfstring{\gu{ભરોસો} (bharoso)}{ભરોસો (bharoso)}]{\gu{ભરોસો} (bharoso) — confidence (trust in someone)}

\label{lesson:GU-C541-bharoso}



\begin{quote}
Before the new one: say the Gujarati for a solemn oath, then the Gujarati for an undertaking.
\end{quote}

\subsection*{You'll want to know: \gu{ભરોસો}}
\textbf{\gu{ભરોસો}} — \emph{bharoso} — \textquotedblleft{}confidence (trust in someone)\textquotedblright{}.

\begin{quote}
\gu{જો} \gu{તમે} \gu{વાયદો} \gu{પાળશો} \gu{તો} \gu{લોકો} \gu{તમારા} \gu{પર} \gu{ભરોસો} \gu{રાખશે}.
\end{quote}

\begin{quote}
\emph{jo tame vāydo pāḷśo to loko tamārā par bharoso rākhśe.}
\end{quote}

\begin{quote}
If you keep your word, people will trust you.
\end{quote}

A masculine noun.

\begin{cousinweb}
\textbf{\gu{ભરોસો}} (\emph{bharoso}) goes back to a reconstructed \emph{*bharavaśya}, \textquotedblleft{}dependence on support\textquotedblright{}, and CDIAL lists the Gujarati word, \textquotedblleft{}confidence\textquotedblright{} (CDIAL 9398). One places it on someone: \textbf{\gu{તમારા} \gu{પર} \gu{ભરોસો}}.
\end{cousinweb}

\begin{grammarlens}[title={What you've built}]
One more word for saying what happens if.
\end{grammarlens}

\subsection*{Your turn}
\begin{itemize}
\raggedright
  \item \emph{Say it:} \emph{bharoso}
  \item \emph{Say it:} \emph{bharoso}, once more
  \item \emph{Say it:} \emph{jo tame vāydo pāḷśo to loko tamārā par bharoso rākhśe.} — If you keep your word, people will trust you.
  \item \emph{Recall:} say the Gujarati for a solemn oath, then the Gujarati for an undertaking, then say \emph{bharoso} again
\end{itemize}

\begin{tcolorbox}[breakable,colback=teal!4,colframe=teal!35!black,title={Before you move on}]
What does \gu{ભરોસો} mean? (\textquotedblleft{}Confidence (trust in someone)\textquotedblright{}.) Say it once more.
\end{tcolorbox}
\section[\texorpdfstring{\gu{છટકબારી} (chaṭakbārī)}{છટકબારી (chaṭakbārī)}]{\gu{છટકબારી} (chaṭakbārī) — a loophole}

\label{lesson:GU-C541-chatakbari}



\begin{quote}
Before the new one: say the Gujarati for an undertaking, then the Gujarati for confidence.
\end{quote}

\subsection*{You'll want to know: \gu{છટકબારી}}
\textbf{\gu{છટકબારી}} — \emph{chaṭakbārī} — \textquotedblleft{}a loophole\textquotedblright{}.

\begin{quote}
\gu{જો} \gu{નિયમમાં} \gu{કોઈ} \gu{છટકબારી} \gu{હશે} \gu{તો} \gu{કોઈ} \gu{તે} \gu{નિયમ} \gu{નહીં} \gu{પાળે}.
\end{quote}

\begin{quote}
\emph{jo niyammā̃ koī chaṭakbārī haśe to koī te niyam nahī̃ pāḷe.}
\end{quote}

\begin{quote}
If there is a loophole in the rule, nobody will keep it.
\end{quote}

A feminine noun.

\begin{cousinweb}
\textbf{\gu{છટકબારી}} (\emph{chaṭakbārī}) is built from \textbf{\gu{છટકવું}}, \textquotedblleft{}to slip away\textquotedblright{}, which CDIAL files under \emph{*chaṭṭ} (CDIAL 4968), and \textbf{\gu{બારી}}, \textquotedblleft{}a window\textquotedblright{}: a window to slip out through.
\end{cousinweb}

\begin{grammarlens}[title={What you've built}]
One more word for saying what happens if.
\end{grammarlens}

\subsection*{Your turn}
\begin{itemize}
\raggedright
  \item \emph{Say it:} \emph{chaṭakbārī}
  \item \emph{Say it:} \emph{chaṭakbārī}, once more
  \item \emph{Say it:} \emph{jo niyammā̃ koī chaṭakbārī haśe to koī te niyam nahī̃ pāḷe.} — If there is a loophole in the rule, nobody will keep it.
  \item \emph{Recall:} say the Gujarati for an undertaking, then the Gujarati for confidence, then say \emph{chaṭakbārī} again
\end{itemize}

\begin{tcolorbox}[breakable,colback=teal!4,colframe=teal!35!black,title={Before you move on}]
What does \gu{છટકબારી} mean? (\textquotedblleft{}A loophole\textquotedblright{}.) Say it once more.
\end{tcolorbox}
\section[\texorpdfstring{\gu{પ્રસ્તાવ} (prastāv)}{પ્રસ્તાવ (prastāv)}]{\gu{પ્રસ્તાવ} (prastāv) — a proposal}

\label{lesson:GU-C541-prastav}



\begin{quote}
Before the new one: say the Gujarati for confidence, then the Gujarati for a loophole.
\end{quote}

\subsection*{You'll want to know: \gu{પ્રસ્તાવ}}
\textbf{\gu{પ્રસ્તાવ}} — \emph{prastāv} — \textquotedblleft{}a proposal\textquotedblright{}.

\begin{quote}
\gu{જો} \gu{તમને} \gu{મારો} \gu{પ્રસ્તાવ} \gu{ગમશે} \gu{તો} \gu{આપણે} \gu{સાથે} \gu{કામ} \gu{કરીશું}.
\end{quote}

\begin{quote}
\emph{jo tamne māro prastāv gamśe to āpaṇe sāthe kām karīśũ.}
\end{quote}

\begin{quote}
If you like my proposal, we'll work together.
\end{quote}

A masculine noun.

\begin{cousinweb}
\textbf{\gu{પ્રસ્તાવ}} (\emph{prastāv}) is a tatsama, taken from Sanskrit \emph{prastāva}, \textquotedblleft{}an introduction, a prelude\textquotedblright{} (Monier-Williams); in Gujarati it is a proposal put forward.
\end{cousinweb}

\begin{grammarlens}[title={What you've built}]
One more word for saying what happens if.
\end{grammarlens}

\subsection*{Your turn}
\begin{itemize}
\raggedright
  \item \emph{Say it:} \emph{prastāv}
  \item \emph{Say it:} \emph{prastāv}, once more
  \item \emph{Say it:} \emph{jo tamne māro prastāv gamśe to āpaṇe sāthe kām karīśũ.} — If you like my proposal, we'll work together.
  \item \emph{Recall:} say the Gujarati for confidence, then the Gujarati for a loophole, then say \emph{prastāv} again
\end{itemize}

\begin{tcolorbox}[breakable,colback=teal!4,colframe=teal!35!black,title={Before you move on}]
What does \gu{પ્રસ્તાવ} mean? (\textquotedblleft{}A proposal\textquotedblright{}.) Say it once more.
\end{tcolorbox}
